\section{Results}
\label{sec:results}

\subsection{Target Organization of Results}

\begin{itemize}[leftmargin=*]
    \item The final paper should present results in research-question order, not in the order of the currently available scripts.
    \item Proposed narrative:
    \begin{itemize}[leftmargin=*]
        \item overall performance on reachability tasks;
        \item effect of structural factors;
        \item control-flow vs data-flow comparison;
        \item reasoning vs non-reasoning comparison and effect of the requested budget.
    \end{itemize}
    \item Current status:
    \begin{itemize}[leftmargin=*]
        \item the tables and figures below mostly document the preliminary control-flow study;
        \item subsections RQ1 to RQ4 currently remain placeholders for writing and future analysis.
    \end{itemize}
\end{itemize}

\begin{figure*}[t]
\centering
\begin{tikzpicture}[rb diagram,
  pipeline block/.style={rb block, text width=2.6cm, minimum height=2.1cm},
  pipeline prompt/.style={rb block, text width=4.1cm, minimum height=1cm},
  pipeline oracle/.style={rb edge, dashed}]

\node[rb heading, anchor=west] at (0,1.4) {(a) Dataset construction};
\node[pipeline block] (settings) at (1.5,0)
  {\textbf{Settings}\\Methods, guards, parameters\\Comments, variables, loops};
\node[pipeline block] (java) at (4.9,0)
  {\textbf{Program generation}\\Method tree\\Java class; fixed execution};
\node[pipeline block] (questions) at (8.3,0)
  {\textbf{Oracle and questions}\\Executable calls\\Source--target pairs};
\node[pipeline block, text width=4.3cm, draw=rbInk, line width=0.95pt,
  fill=modelblue!5] (dataset) at (12.9,0)
  {\textbf{\benchmark: final dataset}\\Java classes and reachability questions\\Boolean labels and executable-graph oracles};
\foreach \from/\to in {settings/java,java/questions,questions/dataset}
  \draw[rb edge] (\from.east) -- (\to.west);

\node[rb heading, anchor=west] at (0,-1.85) {(b) Model evaluation};
\node[pipeline prompt] (co) at (4.6,-2.7)
  {\textbf{CO: Class-only}\\Class; entry context and determinism implicit};
\node[pipeline prompt] (ee) at (4.6,-4.25)
  {\textbf{EE: Explicit execution}\\Class + \texttt{Main}; determinism explicit};
\node[pipeline block, text width=2.35cm] (model) at (9.15,-3.475)
  {\textbf{LLM inference}\\For each variant:\\Boolean answer and trace};
\node[pipeline block, text width=3cm] (scoring) at (13.1,-3.475)
  {\textbf{Independent scoring}\\Final-answer correctness\\Trace correctness};

% Only the class and question enter either model prompt.
\draw[rb edge] (dataset.south west) -- (10.575,-1.25)
  -- (-0.2,-1.25) -- (-0.2,-2.7) -- (co.west);
\node[rb annotation] at (6.05,-1.25) {Same class and question};
\draw[rb edge] (-0.2,-2.7) -- (-0.2,-4.25) -- (ee.west);
\draw[rb edge] (co.east) -- (7.55,-2.7) -- (7.55,-3.475) -- (model.west);
\draw[rb edge, -] (ee.east) -- (7.55,-4.25) -- (7.55,-3.475);
\draw[rb edge] (model.east) -- (scoring.west);

% Labels and executable graphs are reserved for evaluation.
\draw[pipeline oracle] (dataset.south) -- (12.9,-1.55)
  -- (13.1,-1.55) -- (scoring.north);
\node[rb annotation, anchor=west, align=left] at (13.25,-1.75)
  {Shared labels\\and oracles};
\end{tikzpicture}
\caption{Construction and evaluation of \benchmark. The construction algorithm produces the final dataset. CO and EE present the same classes and questions with different execution information. Labels and executable-graph oracles are used only for scoring the returned answer and trace.}
\Description{The construction algorithm generates Java classes, records executable calls and selects reachability questions to form ReachabilityBench. The same class and question feed either a Class-only prompt or an Explicit execution prompt with Main and deterministic-execution instructions. Model answers and traces are scored separately against the shared labels and executable-graph oracles, which are not supplied to the model.}
\label{fig:pipeline}
\end{figure*}


\subsection{\researchq{1}: Overall Performance on Reachability Tasks}

\begin{itemize}[leftmargin=*]
    \item Present here the final high-level view of model performance.
    \begin{itemize}[leftmargin=*]
        \item control-flow / data-flow comparison;
        \item reasoning / non-reasoning comparison;
        \item final metrics and trace metrics.
    \end{itemize}
\end{itemize}

\subsection{\researchq{2}: Effect of Structural Difficulty Factors}

\begin{itemize}[leftmargin=*]
    \item Group together here the effects of depth, context, branching, and distractors.
    \item Distinguish clearly between:
    \begin{itemize}[leftmargin=*]
        \item factors intrinsic to the control skeleton;
        \item factors that determine whether guarded calls execute in the data-flow family.
    \end{itemize}
\end{itemize}

\subsection{\researchq{3}: Comparison Between Control Flow and Data Flow}

\begin{itemize}[leftmargin=*]
    \item Show what changes when the same reachability question does or does not require reasoning over values.
    \item Look for useful contrastive cases:
    \begin{itemize}[leftmargin=*]
        \item control flow reachable but data flow unreachable;
        \item control flow easy but data flow difficult;
        \item effects specific to condition evaluation and guarded calls.
    \end{itemize}
\end{itemize}

\subsection{\researchq{4}: Reasoning Models, Non-Reasoning Models, and Requested Budget}

\begin{itemize}[leftmargin=*]
    \item Compare gaps in level, robustness, and failure types.
    \item Check whether the gains of reasoning models are concentrated mainly:
    \begin{itemize}[leftmargin=*]
        \item on the final answer;
        \item on the structural validity of the trace;
        \item on robustness to difficulty factors.
    \end{itemize}
    \item For reasoning models, compare performance across requested reasoning budgets and identify saturation.
    \item A preliminary result is already visible: the requested reasoning budget helps, but differently across models.
    \item Keep a comparative narrative:
    \begin{itemize}[leftmargin=*]
        \item the strongest model overall;
        \item the model that saturates early;
        \item the model whose final answers improve more than its traces.
    \end{itemize}
    \item Key interpretation: final-answer accuracy alone overestimates reasoning quality.
\end{itemize}

\begin{table*}[t]
\centering
\caption{Condition-aware performance over 46,656 questions per model (\%).}
\label{tab:overall}
\scriptsize
\begin{tabular}{lrrr}
\toprule
Model & Final answer & Path & No response \\
\midrule
Gemma-4-26B-A4B & 84.21 & 55.90 & 2.24 \\
Qwen3.6-35B-A3B & \textbf{87.98} & \textbf{67.43} & \textbf{0.83} \\
\bottomrule
\end{tabular}
\end{table*}


\begin{figure*}[t]
    \centering
    \begin{tikzpicture}[x=1cm,y=1cm]
        \draw[modelblue,line width=1.1pt] (3.75,3.62) -- +(0.38,0); \node[right,font=\scriptsize] at (4.20,3.62) {Gemma};
        \draw[modelorange,line width=1.1pt] (5.60,3.62) -- +(0.38,0); \node[right,font=\scriptsize] at (6.05,3.62) {GPT-OSS};
        \draw[modelgreen,line width=1.1pt] (7.75,3.62) -- +(0.38,0); \node[right,font=\scriptsize] at (8.20,3.62) {Qwen};
        \begin{scope}[xshift=0cm]
        \node[rb panel title] at (2.4,3.15) {Final-answer accuracy};
        \draw[rb axis] (0,0) -- (4.8,0);
        \draw[rb axis] (0,0) -- (0,2.6);
        \draw[rb grid] (0,0.00) -- (4.8,0.00);
        \node[left,font=\scriptsize] at (-0.06,0.00) {60\%};
        \draw[rb grid] (0,0.65) -- (4.8,0.65);
        \node[left,font=\scriptsize] at (-0.06,0.65) {70\%};
        \draw[rb grid] (0,1.30) -- (4.8,1.30);
        \node[left,font=\scriptsize] at (-0.06,1.30) {80\%};
        \draw[rb grid] (0,1.95) -- (4.8,1.95);
        \node[left,font=\scriptsize] at (-0.06,1.95) {90\%};
        \draw[rb grid] (0,2.60) -- (4.8,2.60);
        \node[left,font=\scriptsize] at (-0.06,2.60) {100\%};
        \draw[rb axis] (0.00,0) -- +(0,-0.05);
        \node[below,font=\scriptsize] at (0.00,-0.10) {512};
        \draw[rb axis] (0.96,0) -- +(0,-0.05);
        \node[below,font=\scriptsize] at (0.96,-0.10) {1k};
        \draw[rb axis] (1.92,0) -- +(0,-0.05);
        \node[below,font=\scriptsize] at (1.92,-0.10) {2k};
        \draw[rb axis] (2.88,0) -- +(0,-0.05);
        \node[below,font=\scriptsize] at (2.88,-0.10) {4k};
        \draw[rb axis] (3.84,0) -- +(0,-0.05);
        \node[below,font=\scriptsize] at (3.84,-0.10) {8k};
        \draw[rb axis] (4.80,0) -- +(0,-0.05);
        \node[below,font=\scriptsize] at (4.80,-0.10) {16k};
        \draw[modelblue,line width=1.1pt] (0.000,1.708) -- (0.960,1.926) -- (1.920,1.991) -- (2.880,2.007) -- (3.840,2.058) -- (4.800,2.039);
        \fill[modelblue] (0.000,1.708) circle (1.2pt);
        \fill[modelblue] (0.960,1.926) circle (1.2pt);
        \fill[modelblue] (1.920,1.991) circle (1.2pt);
        \fill[modelblue] (2.880,2.007) circle (1.2pt);
        \fill[modelblue] (3.840,2.058) circle (1.2pt);
        \fill[modelblue] (4.800,2.039) circle (1.2pt);
        \draw[modelorange,line width=1.1pt] (0.000,0.322) -- (0.960,0.710) -- (1.920,1.239) -- (2.880,1.588) -- (3.840,1.751) -- (4.800,1.819);
        \fill[modelorange] (0.000,0.322) circle (1.2pt);
        \fill[modelorange] (0.960,0.710) circle (1.2pt);
        \fill[modelorange] (1.920,1.239) circle (1.2pt);
        \fill[modelorange] (2.880,1.588) circle (1.2pt);
        \fill[modelorange] (3.840,1.751) circle (1.2pt);
        \fill[modelorange] (4.800,1.819) circle (1.2pt);
        \draw[modelgreen,line width=1.1pt] (0.000,0.868) -- (0.960,1.363) -- (1.920,1.825) -- (2.880,2.274) -- (3.840,2.484) -- (4.800,2.531);
        \fill[modelgreen] (0.000,0.868) circle (1.2pt);
        \fill[modelgreen] (0.960,1.363) circle (1.2pt);
        \fill[modelgreen] (1.920,1.825) circle (1.2pt);
        \fill[modelgreen] (2.880,2.274) circle (1.2pt);
        \fill[modelgreen] (3.840,2.484) circle (1.2pt);
        \fill[modelgreen] (4.800,2.531) circle (1.2pt);
        \node[below,font=\scriptsize] at (2.4,-0.48) {Requested budget};
        \end{scope}
        \begin{scope}[xshift=6.45cm]
        \node[rb panel title] at (2.4,3.15) {Exact-chain accuracy};
        \draw[rb axis] (0,0) -- (4.8,0);
        \draw[rb axis] (0,0) -- (0,2.6);
        \draw[rb grid] (0,0.00) -- (4.8,0.00);
        \node[left,font=\scriptsize] at (-0.06,0.00) {20\%};
        \draw[rb grid] (0,0.40) -- (4.8,0.40);
        \node[left,font=\scriptsize] at (-0.06,0.40) {30\%};
        \draw[rb grid] (0,0.80) -- (4.8,0.80);
        \node[left,font=\scriptsize] at (-0.06,0.80) {40\%};
        \draw[rb grid] (0,1.20) -- (4.8,1.20);
        \node[left,font=\scriptsize] at (-0.06,1.20) {50\%};
        \draw[rb grid] (0,1.60) -- (4.8,1.60);
        \node[left,font=\scriptsize] at (-0.06,1.60) {60\%};
        \draw[rb grid] (0,2.00) -- (4.8,2.00);
        \node[left,font=\scriptsize] at (-0.06,2.00) {70\%};
        \draw[rb grid] (0,2.40) -- (4.8,2.40);
        \node[left,font=\scriptsize] at (-0.06,2.40) {80\%};
        \draw[rb axis] (0.00,0) -- +(0,-0.05);
        \node[below,font=\scriptsize] at (0.00,-0.10) {512};
        \draw[rb axis] (0.96,0) -- +(0,-0.05);
        \node[below,font=\scriptsize] at (0.96,-0.10) {1k};
        \draw[rb axis] (1.92,0) -- +(0,-0.05);
        \node[below,font=\scriptsize] at (1.92,-0.10) {2k};
        \draw[rb axis] (2.88,0) -- +(0,-0.05);
        \node[below,font=\scriptsize] at (2.88,-0.10) {4k};
        \draw[rb axis] (3.84,0) -- +(0,-0.05);
        \node[below,font=\scriptsize] at (3.84,-0.10) {8k};
        \draw[rb axis] (4.80,0) -- +(0,-0.05);
        \node[below,font=\scriptsize] at (4.80,-0.10) {16k};
        \draw[modelblue,line width=1.1pt] (0.000,1.323) -- (0.960,1.455) -- (1.920,1.525) -- (2.880,1.531) -- (3.840,1.557) -- (4.800,1.544);
        \fill[modelblue] (0.000,1.323) circle (1.2pt);
        \fill[modelblue] (0.960,1.455) circle (1.2pt);
        \fill[modelblue] (1.920,1.525) circle (1.2pt);
        \fill[modelblue] (2.880,1.531) circle (1.2pt);
        \fill[modelblue] (3.840,1.557) circle (1.2pt);
        \fill[modelblue] (4.800,1.544) circle (1.2pt);
        \draw[modelorange,line width=1.1pt] (0.000,0.057) -- (0.960,0.357) -- (1.920,0.707) -- (2.880,0.860) -- (3.840,0.967) -- (4.800,0.975);
        \fill[modelorange] (0.000,0.057) circle (1.2pt);
        \fill[modelorange] (0.960,0.357) circle (1.2pt);
        \fill[modelorange] (1.920,0.707) circle (1.2pt);
        \fill[modelorange] (2.880,0.860) circle (1.2pt);
        \fill[modelorange] (3.840,0.967) circle (1.2pt);
        \fill[modelorange] (4.800,0.975) circle (1.2pt);
        \draw[modelgreen,line width=1.1pt] (0.000,0.846) -- (0.960,1.313) -- (1.920,1.799) -- (2.880,2.275) -- (3.840,2.446) -- (4.800,2.448);
        \fill[modelgreen] (0.000,0.846) circle (1.2pt);
        \fill[modelgreen] (0.960,1.313) circle (1.2pt);
        \fill[modelgreen] (1.920,1.799) circle (1.2pt);
        \fill[modelgreen] (2.880,2.275) circle (1.2pt);
        \fill[modelgreen] (3.840,2.446) circle (1.2pt);
        \fill[modelgreen] (4.800,2.448) circle (1.2pt);
        \node[below,font=\scriptsize] at (2.4,-0.48) {Requested budget};
        \end{scope}
    \end{tikzpicture}
    \caption{Performance as a function of requested reasoning budget. The two panels separate gains on the final answer from gains on exact chain, which makes saturation and metric-specific behavior visible.}
    \Description{Two LaTeX line charts showing the evolution of final-answer accuracy and exact chain as requested budget increases for Gemma, GPT-OSS, and Qwen.}
    \label{fig:budget-scaling-latex}
\end{figure*}

\begin{table*}[t]
    \centering
    \caption{Effect of requested reasoning budget on the common evaluation slice. Each cell reports weighted accuracy.}
    \label{tab:budget}
    \small
    \begin{tabular}{r rr rr rr}
        \toprule
        & \multicolumn{2}{c}{Gemma} & \multicolumn{2}{c}{GPT-OSS} & \multicolumn{2}{c}{Qwen} \\
        \cmidrule(lr){2-3}\cmidrule(lr){4-5}\cmidrule(lr){6-7}
        Budget & Final & Chain & Final & Chain & Final & Chain \\
        \midrule
        512 & 86.3\% & 53.1\% & 65.0\% & 21.4\% & 73.4\% & 41.1\% \\
        1024 & 89.6\% & 56.4\% & 70.9\% & 28.9\% & 81.0\% & 52.8\% \\
        2048 & 90.6\% & 58.1\% & 79.1\% & 37.7\% & 88.1\% & 65.0\% \\
        4096 & 90.9\% & 58.3\% & 84.4\% & 41.5\% & 95.0\% & 76.9\% \\
        8192 & 91.7\% & 58.9\% & 86.9\% & 44.2\% & 98.2\% & 81.2\% \\
        16384 & 91.4\% & 58.6\% & 88.0\% & 44.4\% & 98.9\% & 81.2\% \\
        \bottomrule
    \end{tabular}
\end{table*}

\begin{table*}[t]
    \centering
    \caption{Marginal gains between consecutive requested reasoning budgets, in percentage points. This table identifies where additional budget stops buying meaningful performance.}
    \label{tab:marginal-budget}
    \small
    \begin{tabular}{lrrrrrr}
        \toprule
        & \multicolumn{2}{c}{Gemma} & \multicolumn{2}{c}{GPT-OSS} & \multicolumn{2}{c}{Qwen} \\
        \cmidrule(lr){2-3}\cmidrule(lr){4-5}\cmidrule(lr){6-7}
        Budget step & Final & Chain & Final & Chain & Final & Chain \\
        \midrule
        512 $\rightarrow$ 1024 & +3.3 & +3.3 & +6.0 & +7.5 & +7.6 & +11.7 \\
        1024 $\rightarrow$ 2048 & +1.0 & +1.7 & +8.2 & +8.8 & +7.1 & +12.2 \\
        2048 $\rightarrow$ 4096 & +0.3 & +0.2 & +5.4 & +3.8 & +6.9 & +11.9 \\
        4096 $\rightarrow$ 8192 & +0.8 & +0.6 & +2.5 & +2.7 & +3.2 & +4.3 \\
        8192 $\rightarrow$ 16384 & -0.3 & -0.3 & +1.1 & +0.2 & +0.7 & +0.0 \\
        \bottomrule
    \end{tabular}
\end{table*}

\begin{table*}[t]
    \centering
    \caption{Availability of tokenizer-counted emitted reasoning tokens in the recomputed CSV summaries. The value counts generated reasoning text extracted between \texttt{R:} and \texttt{A:} with each model's local tokenizer; it is a proxy for emitted reasoning length, not a runtime usage counter.}
    \label{tab:budget-observability}
    \small
    \begin{tabular}{lrrl}
        \toprule
        Model & CSV rows & Proxy rows & Status \\
        \midrule
        Gemma & 450 & 450 & tokenizer-counted \\
        GPT-OSS & 450 & 450 & tokenizer-counted \\
        Qwen & 450 & 450 & tokenizer-counted \\
        \bottomrule
    \end{tabular}
\end{table*}

\begin{table*}[t]
    \centering
    \caption{Tokenizer-counted emitted reasoning tokens by requested budget. Cells report weighted means on the shared slice. The proxy can slightly exceed the requested budget at 512, because it is computed post hoc from \texttt{R:} text extracted with the local tokenizer, whereas the decoding budget is enforced by the serving stack.}
    \label{tab:used-budget-by-requested}
    \small
    \begin{tabular}{lrrrrrr}
        \toprule
        Model & 512 & 1024 & 2048 & 4096 & 8192 & 16384 \\
        \midrule
        Gemma & 514 & 988 & 1790 & 2855 & 3835 & 4468 \\
        GPT-OSS & 503 & 969 & 1737 & 2710 & 3512 & 3870 \\
        Qwen & 516 & 1022 & 1924 & 3137 & 3871 & 3941 \\
        \bottomrule
    \end{tabular}
\end{table*}

\begin{itemize}[leftmargin=*]
    \item Add a short reading of saturation below the figure.
    \item Keep the caveat that the requested reasoning budget is an interface control, not an internal cost.
\end{itemize}

\begin{figure*}[t]
    \centering
    \begin{tikzpicture}[x=0.72cm,y=1cm]
        \begin{scope}[shift={(0,0)}]
        \node[rb panel title] at (2.75,3.13) {Context 100};
        \draw[rb axis] (0,0) -- (5.5,0);
        \draw[rb axis] (0,0) -- (0,2.6);
        \draw[rb grid] (0,0.00) -- (5.5,0.00);
        \node[left,font=\scriptsize] at (-0.05,0.00) {60\%};
        \draw[rb grid] (0,0.65) -- (5.5,0.65);
        \node[left,font=\scriptsize] at (-0.05,0.65) {70\%};
        \draw[rb grid] (0,1.30) -- (5.5,1.30);
        \node[left,font=\scriptsize] at (-0.05,1.30) {80\%};
        \draw[rb grid] (0,1.95) -- (5.5,1.95);
        \node[left,font=\scriptsize] at (-0.05,1.95) {90\%};
        \draw[rb grid] (0,2.60) -- (5.5,2.60);
        \node[left,font=\scriptsize] at (-0.05,2.60) {100\%};
        \draw[rb axis] (0.00,0) -- +(0,-0.05);
        \node[below,font=\scriptsize] at (0.00,-0.10) {1};
        \draw[rb axis] (0.92,0) -- +(0,-0.05);
        \node[below,font=\scriptsize] at (0.92,-0.10) {5};
        \draw[rb axis] (2.06,0) -- +(0,-0.05);
        \node[below,font=\scriptsize] at (2.06,-0.10) {10};
        \draw[rb axis] (3.21,0) -- +(0,-0.05);
        \node[below,font=\scriptsize] at (3.21,-0.10) {15};
        \draw[rb axis] (4.35,0) -- +(0,-0.05);
        \node[below,font=\scriptsize] at (4.35,-0.10) {20};
        \draw[rb axis] (5.50,0) -- +(0,-0.05);
        \node[below,font=\scriptsize] at (5.50,-0.10) {25};
        \draw[modelblue,line width=0.95pt] (0.000,2.600) -- (0.229,2.519) -- (0.458,2.519) -- (0.688,2.519) -- (0.917,2.519) -- (1.146,2.356) -- (1.375,2.438) -- (1.604,2.356) -- (1.833,2.356) -- (2.062,2.356) -- (2.292,2.519) -- (2.521,2.519) -- (2.750,2.356) -- (2.979,2.275) -- (3.208,2.275) -- (3.438,2.600) -- (3.667,2.194) -- (3.896,2.275) -- (4.125,2.275) -- (4.354,2.519) -- (4.583,2.438) -- (4.812,2.519) -- (5.042,2.519) -- (5.271,2.275) -- (5.500,2.519);
        \draw[modelorange,line width=0.95pt] (0.000,2.438) -- (0.229,2.600) -- (0.458,2.438) -- (0.688,2.438) -- (0.917,2.438) -- (1.146,2.356) -- (1.375,2.519) -- (1.604,2.519) -- (1.833,2.519) -- (2.062,2.356) -- (2.292,2.438) -- (2.521,2.113) -- (2.750,2.275) -- (2.979,1.950) -- (3.208,1.788) -- (3.438,2.438) -- (3.667,2.275) -- (3.896,2.438) -- (4.125,2.113) -- (4.354,2.194) -- (4.583,1.950) -- (4.812,2.356) -- (5.042,1.950) -- (5.271,1.544) -- (5.500,2.194);
        \draw[modelgreen,line width=0.95pt] (0.000,2.600) -- (0.229,2.600) -- (0.458,2.600) -- (0.688,2.600) -- (0.917,2.600) -- (1.146,2.600) -- (1.375,2.519) -- (1.604,2.600) -- (1.833,2.519) -- (2.062,2.600) -- (2.292,2.600) -- (2.521,2.600) -- (2.750,2.600) -- (2.979,2.600) -- (3.208,2.600) -- (3.438,2.600) -- (3.667,2.600) -- (3.896,2.600) -- (4.125,2.600) -- (4.354,2.600) -- (4.583,2.519) -- (4.812,2.600) -- (5.042,2.519) -- (5.271,2.438) -- (5.500,2.600);
        \node[below,font=\scriptsize] at (2.75,-0.48) {Depth};
        \end{scope}
        \begin{scope}[shift={(7.2,0)}]
        \node[rb panel title] at (2.75,3.13) {Context 250};
        \draw[rb axis] (0,0) -- (5.5,0);
        \draw[rb axis] (0,0) -- (0,2.6);
        \draw[rb grid] (0,0.00) -- (5.5,0.00);
        \node[left,font=\scriptsize] at (-0.05,0.00) {60\%};
        \draw[rb grid] (0,0.65) -- (5.5,0.65);
        \node[left,font=\scriptsize] at (-0.05,0.65) {70\%};
        \draw[rb grid] (0,1.30) -- (5.5,1.30);
        \node[left,font=\scriptsize] at (-0.05,1.30) {80\%};
        \draw[rb grid] (0,1.95) -- (5.5,1.95);
        \node[left,font=\scriptsize] at (-0.05,1.95) {90\%};
        \draw[rb grid] (0,2.60) -- (5.5,2.60);
        \node[left,font=\scriptsize] at (-0.05,2.60) {100\%};
        \draw[rb axis] (0.00,0) -- +(0,-0.05);
        \node[below,font=\scriptsize] at (0.00,-0.10) {1};
        \draw[rb axis] (0.92,0) -- +(0,-0.05);
        \node[below,font=\scriptsize] at (0.92,-0.10) {5};
        \draw[rb axis] (2.06,0) -- +(0,-0.05);
        \node[below,font=\scriptsize] at (2.06,-0.10) {10};
        \draw[rb axis] (3.21,0) -- +(0,-0.05);
        \node[below,font=\scriptsize] at (3.21,-0.10) {15};
        \draw[rb axis] (4.35,0) -- +(0,-0.05);
        \node[below,font=\scriptsize] at (4.35,-0.10) {20};
        \draw[rb axis] (5.50,0) -- +(0,-0.05);
        \node[below,font=\scriptsize] at (5.50,-0.10) {25};
        \draw[modelblue,line width=0.95pt] (0.000,2.600) -- (0.229,2.275) -- (0.458,2.356) -- (0.688,2.438) -- (0.917,2.113) -- (1.146,1.950) -- (1.375,2.519) -- (1.604,2.519) -- (1.833,2.194) -- (2.062,2.519) -- (2.292,1.950) -- (2.521,2.275) -- (2.750,2.031) -- (2.979,2.194) -- (3.208,2.194) -- (3.438,2.113) -- (3.667,2.031) -- (3.896,1.950) -- (4.125,2.194) -- (4.354,1.869) -- (4.583,1.950) -- (4.812,2.031) -- (5.042,1.706) -- (5.271,1.706) -- (5.500,1.869);
        \draw[modelorange,line width=0.95pt] (0.000,2.519) -- (0.229,2.275) -- (0.458,2.356) -- (0.688,2.356) -- (0.917,2.356) -- (1.146,1.869) -- (1.375,1.869) -- (1.604,2.031) -- (1.833,1.950) -- (2.062,1.625) -- (2.292,1.788) -- (2.521,1.625) -- (2.750,1.625) -- (2.979,1.788) -- (3.208,1.788) -- (3.438,2.031) -- (3.667,1.544) -- (3.896,1.381) -- (4.125,1.706) -- (4.354,1.544) -- (4.583,1.381) -- (4.812,1.463) -- (5.042,1.788) -- (5.271,1.869) -- (5.500,2.031);
        \draw[modelgreen,line width=0.95pt] (0.000,2.600) -- (0.229,2.600) -- (0.458,2.600) -- (0.688,2.600) -- (0.917,2.600) -- (1.146,2.519) -- (1.375,2.600) -- (1.604,2.600) -- (1.833,2.600) -- (2.062,2.600) -- (2.292,2.438) -- (2.521,2.600) -- (2.750,2.438) -- (2.979,2.519) -- (3.208,2.519) -- (3.438,2.519) -- (3.667,2.600) -- (3.896,2.519) -- (4.125,2.519) -- (4.354,2.600) -- (4.583,2.356) -- (4.812,2.519) -- (5.042,2.438) -- (5.271,2.275) -- (5.500,2.519);
        \node[below,font=\scriptsize] at (2.75,-0.48) {Depth};
        \end{scope}
        \begin{scope}[shift={(14.4,0)}]
        \node[rb panel title] at (2.75,3.13) {Context 500};
        \draw[rb axis] (0,0) -- (5.5,0);
        \draw[rb axis] (0,0) -- (0,2.6);
        \draw[rb grid] (0,0.00) -- (5.5,0.00);
        \node[left,font=\scriptsize] at (-0.05,0.00) {60\%};
        \draw[rb grid] (0,0.65) -- (5.5,0.65);
        \node[left,font=\scriptsize] at (-0.05,0.65) {70\%};
        \draw[rb grid] (0,1.30) -- (5.5,1.30);
        \node[left,font=\scriptsize] at (-0.05,1.30) {80\%};
        \draw[rb grid] (0,1.95) -- (5.5,1.95);
        \node[left,font=\scriptsize] at (-0.05,1.95) {90\%};
        \draw[rb grid] (0,2.60) -- (5.5,2.60);
        \node[left,font=\scriptsize] at (-0.05,2.60) {100\%};
        \draw[rb axis] (0.00,0) -- +(0,-0.05);
        \node[below,font=\scriptsize] at (0.00,-0.10) {1};
        \draw[rb axis] (0.92,0) -- +(0,-0.05);
        \node[below,font=\scriptsize] at (0.92,-0.10) {5};
        \draw[rb axis] (2.06,0) -- +(0,-0.05);
        \node[below,font=\scriptsize] at (2.06,-0.10) {10};
        \draw[rb axis] (3.21,0) -- +(0,-0.05);
        \node[below,font=\scriptsize] at (3.21,-0.10) {15};
        \draw[rb axis] (4.35,0) -- +(0,-0.05);
        \node[below,font=\scriptsize] at (4.35,-0.10) {20};
        \draw[rb axis] (5.50,0) -- +(0,-0.05);
        \node[below,font=\scriptsize] at (5.50,-0.10) {25};
        \draw[modelblue,line width=0.95pt] (0.000,2.113) -- (0.229,2.031) -- (0.458,1.950) -- (0.688,1.869) -- (0.917,1.544) -- (1.146,1.625) -- (1.375,1.625) -- (1.604,1.950) -- (1.833,1.869) -- (2.062,1.788) -- (2.292,1.788) -- (2.521,1.463) -- (2.750,1.625) -- (2.979,1.300) -- (3.208,1.544) -- (3.438,1.300) -- (3.667,1.706) -- (3.896,1.219) -- (4.125,1.544) -- (4.354,1.056) -- (4.583,0.812) -- (4.812,1.381) -- (5.042,1.625) -- (5.271,0.975) -- (5.500,1.056);
        \draw[modelorange,line width=0.95pt] (0.000,2.275) -- (0.229,1.625) -- (0.458,1.788) -- (0.688,1.706) -- (0.917,1.950) -- (1.146,1.463) -- (1.375,1.544) -- (1.604,1.219) -- (1.833,2.031) -- (2.062,1.544) -- (2.292,1.544) -- (2.521,1.056) -- (2.750,1.300) -- (2.979,1.056) -- (3.208,1.137) -- (3.438,1.381) -- (3.667,0.731) -- (3.896,1.463) -- (4.125,1.463) -- (4.354,1.056) -- (4.583,0.812) -- (4.812,1.137) -- (5.042,0.812) -- (5.271,0.650) -- (5.500,0.488);
        \draw[modelgreen,line width=0.95pt] (0.000,2.600) -- (0.229,2.600) -- (0.458,2.600) -- (0.688,2.519) -- (0.917,2.519) -- (1.146,2.519) -- (1.375,2.600) -- (1.604,2.438) -- (1.833,2.600) -- (2.062,2.356) -- (2.292,2.519) -- (2.521,2.600) -- (2.750,2.519) -- (2.979,2.519) -- (3.208,2.438) -- (3.438,2.275) -- (3.667,2.438) -- (3.896,2.519) -- (4.125,2.275) -- (4.354,2.438) -- (4.583,2.356) -- (4.812,2.519) -- (5.042,2.519) -- (5.271,2.438) -- (5.500,2.275);
        \node[below,font=\scriptsize] at (2.75,-0.48) {Depth};
        \end{scope}
        \draw[modelblue,line width=1.1pt] (7.1,3.62) -- +(0.45,0); \node[right,font=\scriptsize] at (7.6,3.62) {Gemma};
        \draw[modelorange,line width=1.1pt] (9.6,3.62) -- +(0.45,0); \node[right,font=\scriptsize] at (10.1,3.62) {GPT-OSS};
        \draw[modelgreen,line width=1.1pt] (12.7,3.62) -- +(0.45,0); \node[right,font=\scriptsize] at (13.2,3.62) {Qwen};
    \end{tikzpicture}
    \caption{Final-answer accuracy at the maximum requested budget, broken down by context and target depth. This directly tests whether 16\,384 requested tokens are enough to remove sensitivity to depth and context.}
    \Description{Three LaTeX line charts showing final-answer accuracy by depth at the 16\,384 requested budget, with one panel per context size.}
    \label{fig:depth-context-latex}
\end{figure*}

\begin{table*}[t]
    \centering
    \caption{Performance at the maximum requested budget, 16,384 tokens, grouped by depth band and context size. Values are weighted percentages. This is the key diagnostic table for asking what remains hard when the budget is maximal.}
    \label{tab:high-budget-16384}
    \scriptsize
    \begin{tabular}{llrrrrrrrrr}
        \toprule
        Metric & Model & \multicolumn{3}{c}{Depth 1--5} & \multicolumn{3}{c}{Depth 6--15} & \multicolumn{3}{c}{Depth 16--25} \\
        \cmidrule(lr){3-5}\cmidrule(lr){6-8}\cmidrule(lr){9-11}
        & & C100 & C250 & C500 & C100 & C250 & C500 & C100 & C250 & C500 \\
        \midrule
        Final & Gemma & 99.0 & 96.2 & 89.2 & 96.6 & 94.4 & 85.5 & 97.1 & 89.9 & 79.5 \\
        Final & GPT-OSS & 98.0 & 96.5 & 88.8 & 95.1 & 87.6 & 81.4 & 93.0 & 85.8 & 75.4 \\
        Final & Qwen & 100.0 & 100.0 & 99.5 & 99.8 & 99.1 & 98.6 & 99.5 & 98.2 & 97.0 \\
        \midrule
        Chain & Gemma & 83.2 & 76.8 & 64.0 & 70.5 & 62.4 & 47.2 & 64.2 & 49.6 & 33.5 \\
        Chain & GPT-OSS & 73.2 & 72.5 & 58.8 & 51.6 & 42.1 & 35.4 & 42.0 & 35.9 & 23.5 \\
        Chain & Qwen & 95.8 & 97.2 & 95.8 & 83.1 & 85.4 & 83.4 & 71.4 & 71.9 & 69.5 \\
        \midrule
        Trace & Gemma & 64.0 & 59.2 & 54.5 & 27.0 & 22.2 & 13.1 & 16.4 & 10.4 & 4.1 \\
        Trace & GPT-OSS & 60.2 & 61.0 & 51.2 & 16.9 & 11.2 & 9.4 & 8.0 & 4.4 & 1.2 \\
        Trace & Qwen & 73.8 & 77.2 & 80.2 & 34.6 & 38.5 & 37.5 & 21.6 & 23.8 & 23.1 \\
        \bottomrule
    \end{tabular}
\end{table*}

\begin{figure*}[t]
    \centering
    \begin{tikzpicture}[x=0.21cm,y=0.43cm]
        \foreach \x/\lab in {0/0\%,10/10\%,20/20\%,30/30\%,40/40\%} {
            \draw[gray!18] (\x,0) -- (\x,24.2);
            \node[below,font=\scriptsize] at (\x,-0.15) {\lab};
        }
        \draw[gray!45] (0,0) -- (42,0);

        \node[anchor=east,font=\scriptsize] at (-0.7,22.95) {Shorter trace};
        \fill[modelblue] (0,23.40) rectangle (31.7,23.76);
        \node[right,font=\tiny] at (32.0,23.58) {31.7\%};
        \fill[modelorange] (0,22.85) rectangle (38.5,23.21);
        \node[right,font=\tiny] at (38.8,23.03) {38.5\%};
        \fill[modelgreen] (0,22.30) rectangle (18.0,22.66);
        \node[right,font=\tiny] at (18.3,22.48) {18.0\%};

        \node[anchor=east,font=\scriptsize] at (-0.7,19.95) {Over-exploration};
        \fill[modelblue] (0,20.40) rectangle (14.6,20.76);
        \node[right,font=\tiny] at (14.9,20.58) {14.6\%};
        \fill[modelorange] (0,19.85) rectangle (31.5,20.21);
        \node[right,font=\tiny] at (31.8,20.03) {31.5\%};
        \fill[modelgreen] (0,19.30) rectangle (7.5,19.66);
        \node[right,font=\tiny] at (7.8,19.48) {7.5\%};

        \node[anchor=east,font=\scriptsize] at (-0.7,16.95) {Alternative branch};
        \fill[modelblue] (0,17.40) rectangle (12.4,17.76);
        \node[right,font=\tiny] at (12.7,17.58) {12.4\%};
        \fill[modelorange] (0,16.85) rectangle (26.6,17.21);
        \node[right,font=\tiny] at (26.9,17.03) {26.6\%};
        \fill[modelgreen] (0,16.30) rectangle (6.2,16.66);
        \node[right,font=\tiny] at (6.5,16.48) {6.2\%};

        \node[anchor=east,font=\scriptsize] at (-0.7,13.95) {Invalid transition};
        \fill[modelblue] (0,14.40) rectangle (11.8,14.76);
        \node[right,font=\tiny] at (12.1,14.58) {11.8\%};
        \fill[modelorange] (0,13.85) rectangle (22.7,14.21);
        \node[right,font=\tiny] at (23.0,14.03) {22.7\%};
        \fill[modelgreen] (0,13.30) rectangle (4.9,13.66);
        \node[right,font=\tiny] at (5.2,13.48) {4.9\%};

        \node[anchor=east,font=\scriptsize] at (-0.7,10.95) {Missed reachable target};
        \fill[modelblue] (0,11.40) rectangle (8.1,11.76);
        \node[right,font=\tiny] at (8.4,11.58) {8.1\%};
        \fill[modelorange] (0,10.85) rectangle (18.3,11.21);
        \node[right,font=\tiny] at (18.6,11.03) {18.3\%};
        \fill[modelgreen] (0,10.30) rectangle (7.0,10.66);
        \node[right,font=\tiny] at (7.3,10.48) {7.0\%};

        \node[anchor=east,font=\scriptsize] at (-0.7,7.95) {Longer trace};
        \fill[modelblue] (0,8.40) rectangle (7.8,8.76);
        \node[right,font=\tiny] at (8.1,8.58) {7.8\%};
        \fill[modelorange] (0,7.85) rectangle (19.3,8.21);
        \node[right,font=\tiny] at (19.6,8.03) {19.3\%};
        \fill[modelgreen] (0,7.30) rectangle (4.9,7.66);
        \node[right,font=\tiny] at (5.2,7.48) {4.9\%};

        \node[anchor=east,font=\scriptsize] at (-0.7,4.95) {Ordering issue};
        \fill[modelblue] (0,5.40) rectangle (6.2,5.76);
        \node[right,font=\tiny] at (6.5,5.58) {6.2\%};
        \fill[modelorange] (0,4.85) rectangle (9.5,5.21);
        \node[right,font=\tiny] at (9.8,5.03) {9.5\%};
        \fill[modelgreen] (0,4.30) rectangle (7.9,4.66);
        \node[right,font=\tiny] at (8.2,4.48) {7.9\%};

        \node[anchor=east,font=\scriptsize] at (-0.7,1.95) {Empty prediction};
        \fill[modelblue] (0,2.40) rectangle (4.0,2.76);
        \node[right,font=\tiny] at (4.3,2.58) {4.0\%};
        \fill[modelorange] (0,1.85) rectangle (6.7,2.21);
        \node[right,font=\tiny] at (7.0,2.03) {6.7\%};
        \fill[modelgreen] (0,1.30) rectangle (6.1,1.66);
        \node[right,font=\tiny] at (6.4,1.48) {6.1\%};

        \fill[modelblue] (9.5,25.05) rectangle ++(0.9,0.30);
        \node[right,font=\scriptsize] at (10.7,25.20) {Gemma};
        \fill[modelorange] (17.0,25.05) rectangle ++(0.9,0.30);
        \node[right,font=\scriptsize] at (18.2,25.20) {GPT-OSS};
        \fill[modelgreen] (27.0,25.05) rectangle ++(0.9,0.30);
        \node[right,font=\scriptsize] at (28.2,25.20) {Qwen};
        \node[below,font=\scriptsize] at (21,-0.80) {Share of all cases};
    \end{tikzpicture}
    \caption{Dominant failure families from the in-depth analysis for chain-level errors. Families are normalized by all evaluated cases and are not mutually exclusive, because a single trace can expose several local characteristics.}
    \Description{Grouped horizontal LaTeX bar chart of aggregated failure families for the three models.}
    \label{fig:error-taxonomy-latex}
\end{figure*}

\begin{table*}[!t]
\centering
\caption{Planned error-family taxonomy. Counts remain TODOs until final reruns and dashboard inspection are complete.}
\label{tab:error-taxonomy}
\tiny
\setlength{\tabcolsep}{2pt}
\begin{tabularx}{\textwidth}{p{0.15\textwidth}p{0.27\textwidth}p{0.21\textwidth}p{0.055\textwidth}p{0.055\textwidth}p{0.15\textwidth}}
\toprule
Error family & Definition & Detection method & $n$ & \% & Example \\
\midrule
Syntactic-call bias & Follows a visible call that is not executable under generated guards. & Compare predicted edges with executable graph. & \TODO{n} & \TODO{\%} & false guarded call \\
Missing intermediate nodes & Skips one or more oracle methods. & Trace alignment. & \TODO{n} & \TODO{\%} & $A,C$ for $A,B,C$ \\
Early stopping & Ends before oracle exploration is complete. & Prefix and coverage checks. & \TODO{n} & \TODO{\%} & source only \\
Over-exploration & Adds reachable but non-oracle exploration after the stopping point. & Oracle stop-condition check. & \TODO{n} & \TODO{\%} & past target \\
Invalid transition & Predicts an edge absent from the executable graph. & Edge-validity check. & \TODO{n} & \TODO{\%} & cross branch \\
Target hallucination & Mentions a target or method name not present in the generated class. & Method-list check. & \TODO{n} & \TODO{\%} & unknown name \\
Positive shortcut & Gives a reachable label with a shortest-path-like trace that omits required exploration. & Canonical-trace comparison. & \TODO{n} & \TODO{\%} & path only \\
Negative over-exploration & Negative trace explores outside the source-reachable region. & Reachable-region check. & \TODO{n} & \TODO{\%} & sibling subtree \\
Wrong source & Starts the trace from a method other than the queried source. & First-node check. & \TODO{n} & \TODO{\%} & target first \\
Repeated methods & Repeats methods in an acyclic tree trace. & Duplicate check. & \TODO{n} & \TODO{\%} & $A,B,A$ \\
Order violation & Uses valid methods but not in canonical order. & Ordered oracle comparison. & \TODO{n} & \TODO{\%} & sibling swap \\
Invalid JSON / schema failure & Output cannot be parsed into the required schema. & Parser status. & \TODO{n} & \TODO{\%} & missing field \\
\bottomrule
\end{tabularx}
\end{table*}

\begin{table*}[t]
    \centering
    \caption{Evolution of dominant chain-level failure families across requested reasoning budgets. Rates are percentages of all cases for the corresponding model and budget. Families are derived from deeper-analysis tags: shorter traces use \texttt{length\_shorter}, over-exploration uses \texttt{extra\_known\_internal} or \texttt{extra\_known\_leaf}, alternative branches use \texttt{coverage\_overlap}, invalid transitions use \texttt{structure\_with\_jump}, and missed reachable targets use \texttt{target\_expected\_missing}.}
    \label{tab:failure-budget-evolution}
    \scriptsize
    \begin{tabular}{llrrrrrr}
        \toprule
        Model & Failure family & 512 & 1k & 2k & 4k & 8k & 16k \\
        \midrule
        Gemma & Shorter trace & 32.2 & 31.9 & 30.9 & 31.7 & 31.1 & 32.3 \\
        Gemma & Over-exploration & 13.6 & 14.0 & 14.1 & 14.8 & 15.3 & 14.9 \\
        Gemma & Alternative branch & 12.1 & 12.1 & 11.8 & 12.7 & 12.9 & 12.8 \\
        Gemma & Invalid transition & 10.1 & 10.7 & 11.4 & 12.3 & 12.9 & 12.5 \\
        Gemma & Reachable target missed & 9.8 & 8.6 & 8.0 & 7.9 & 7.1 & 7.5 \\
        \midrule
        GPT-OSS & Shorter trace & 39.8 & 37.7 & 35.1 & 38.0 & 39.8 & 40.9 \\
        GPT-OSS & Over-exploration & 37.1 & 33.1 & 27.8 & 22.9 & 22.1 & 21.7 \\
        GPT-OSS & Alternative branch & 42.4 & 35.5 & 26.6 & 20.3 & 18.1 & 16.5 \\
        GPT-OSS & Invalid transition & 27.9 & 23.1 & 18.6 & 15.8 & 13.8 & 13.0 \\
        GPT-OSS & Reachable target missed & 28.8 & 24.9 & 18.5 & 14.1 & 12.1 & 11.2 \\
        \midrule
        Qwen & Shorter trace & 29.6 & 25.7 & 18.6 & 12.2 & 10.4 & 11.4 \\
        Qwen & Over-exploration & 15.7 & 10.5 & 5.9 & 3.1 & 1.5 & 1.4 \\
        Qwen & Alternative branch & 14.8 & 10.2 & 6.0 & 3.0 & 1.5 & 1.5 \\
        Qwen & Invalid transition & 9.2 & 5.2 & 3.9 & 2.4 & 1.1 & 1.2 \\
        Qwen & Reachable target missed & 16.6 & 12.3 & 7.1 & 3.4 & 1.5 & 1.0 \\
        \bottomrule
    \end{tabular}
\end{table*}

\begin{table*}[t]
    \centering
    \caption{Representative local motifs for the deeper-analysis failure families. Notation: \texttt{S}=source, \texttt{T}=target, \texttt{B}=node present in both reference and prediction, \texttt{E}=reference-only node, \texttt{P}=prediction-only node. Motifs compact recurring local-graph signatures rather than individual examples.}
    \label{tab:local-motif-examples}
    \small
    \begin{tabularx}{\textwidth}{p{0.20\textwidth}p{0.25\textwidth}X}
        \toprule
        Failure family & Example local motif & Interpretation \\
        \midrule
        Shorter trace & \texttt{S -> B -> B; B -> E -> ...} & The prediction follows a valid prefix and then stops before the reference suffix. \\
        Over-exploration & \texttt{S -> B -> P -> P} & The model keeps traversing prediction-only nodes after leaving the reference path. \\
        Alternative branch & \texttt{B -> E; B -> P} & The prediction shares a prefix, then chooses a sibling or nearby branch instead of the oracle branch. \\
        Invalid transition & \texttt{B -> P} with no graph edge & The predicted transition is not a valid local call edge. \\
        Reachable target missed & \texttt{S -> B -> E -> T} & The target is reachable in the oracle trace, but the predicted trace omits the target-bearing suffix. \\
        Longer trace & \texttt{S -> B -> T -> P} & The prediction reaches the relevant region but continues beyond the reference trace. \\
        Order issue & \texttt{S -> B1, S -> B2} with swapped visits & The local calls are plausible, but the traversal order disagrees with the reference order. \\
        Empty prediction & \texttt{S -> E, S -> E} & The reference has reachable outgoing structure, but the extracted prediction is empty. \\
        Hallucinated target & \texttt{B -> P -> T} & The prediction reaches the target through a prediction-only suffix absent from the oracle path. \\
        Target at wrong rank & \texttt{S -> T -> B} vs. \texttt{S -> B -> T} & The target appears, but at the wrong position in the traversal. \\
        Repeated methods & \texttt{B -> P -> P} & A method or occurrence is repeated in the prediction without a matching repetition in the oracle trace. \\
        Unresolved node & \texttt{B -> P(unresolved)} & The prediction introduces a method name that cannot be resolved in the parsed call graph. \\
        Wrong source & \texttt{P -> P; S -> E} & The predicted trace starts from a node other than the required source anchor. \\
        \bottomrule
    \end{tabularx}
\end{table*}

\begin{table*}[t]
    \centering
    \caption{Example of a correct final answer with an incorrect trace. The target is reachable, but the model skips the intermediate method, so final-answer accuracy overestimates path fidelity.}
    \label{tab:example}
    \small
    \begin{tabularx}{\textwidth}{p{0.17\textwidth}X}
        \toprule
        Field & Value \\
        \midrule
        Question & Does `loadValueFromDatabase` call `calculateMetricInThread`, directly or indirectly? \\
        Budget / depth & 512 tokens / depth 2 \\
        Oracle chain & \texttt{loadValueFromDatabase} \(\rightarrow\) \texttt{dispatchEventForExecution} \(\rightarrow\) \texttt{calculateMetricInThread} \\
        Predicted trace & \texttt{loadValueFromDatabase} \(\rightarrow\) \texttt{calculateMetricInThread} \\
        Interpretation & The final answer is \true, but the reasoning trace contains a shortcut edge that is not present in the oracle call path. \\
        \bottomrule
    \end{tabularx}
\end{table*}

\begin{table*}[t]
    \centering
    \caption{Preliminary contrasts between factors on the shared slice. Values are deltas in percentage points between stressed and baseline conditions, computed by comparing the highest indicated distractor level to the lowest while aggregating over the shared model/context/depth/budget slice. These are descriptive contrasts, not adjusted causal estimates.}
    \label{tab:factors}
    \small
    \begin{tabular}{l rr rr rr}
        \toprule
        & \multicolumn{2}{c}{Gemma} & \multicolumn{2}{c}{GPT-OSS} & \multicolumn{2}{c}{Qwen} \\
        \cmidrule(lr){2-3}\cmidrule(lr){4-5}\cmidrule(lr){6-7}
        Factor contrast & Final & Chain & Final & Chain & Final & Chain \\
        \midrule
        comment: 0 -> 1 & -6.7 pp & -10.5 pp & -6.4 pp & -7.2 pp & -2.2 pp & -4.0 pp \\
        variable: 0 -> 2 & -5.1 pp & -8.1 pp & -4.5 pp & -5.0 pp & -2.4 pp & -3.9 pp \\
        loop: 0 -> 3 & -3.6 pp & -8.5 pp & -7.6 pp & -8.2 pp & -2.0 pp & -5.0 pp \\
        if: 1 -> 2 & -3.6 pp & -9.5 pp & -7.5 pp & -8.3 pp & -2.4 pp & -6.1 pp \\
        parameter: 0 -> 2 & -2.4 pp & -3.2 pp & -1.9 pp & -2.4 pp & -0.8 pp & +0.8 pp \\
        \bottomrule
    \end{tabular}
\end{table*}


\begin{itemize}[leftmargin=*]
    \item Use these figures and tables to support the analysis of trace fidelity, failure modes, and budget sensitivity in \researchq{4}.
\end{itemize}
